% explainer-source-hash: 57bcbdec561bab334481f03dc7f1b4d5b8f071f1c30602f865e48dc1657f0e38
% explainer-source-commit: 40e6ab02e9e59f22535727fa3f73442500a09963
% explainer-generated: 2026-07-16
% Regenerate with /explain-all (see WIRING.md). Do not edit this stamp by hand:
% it is how the toolchain knows whether this document still matches the code.
\documentclass[11pt,a4paper]{article}
\input{preamble}

\title{\textbf{Hi-Tec Shopify Automation}\\[0.3em]\large A Brief}
\author{Portfolio explainer, built from the repository source}
\date{}

\begin{document}
\maketitle
\thispagestyle{empty}

\section{What it is}

Three small Python programs, written solo as freelance work for Hi-Tec Bearings in
2023. The client sells bearings through a Shopify shop with thousands of product
variants. The stock figures they trust live in a spreadsheet exported from their
warehouse system; what the shop tells customers lives in a different spreadsheet
exported from Shopify. Reconciling the two by hand took hours per cycle. These
tools do it automatically.

\begin{keyidea}{The defining constraint: no API}
Nothing here talks to Shopify over the network. No API client, no access token, no
shop domain. A human exports a spreadsheet from the Shopify admin, runs a tool,
reviews the output, and re-imports it. Everything in the design follows from that.
\end{keyidea}

\begin{figure}[H]
\centering
\resizebox{\textwidth}{!}{
\begin{tikzpicture}[node distance=7mm and 11mm]
  \node[extbox] (store) {Shopify store};
  \node[box, right=of store] (exp) {human exports};
  \node[store, right=of exp] (csv) {export\\CSV / XLSX};
  \node[box, right=of csv] (tool) {tool runs\\(offline)};
  \node[store, right=of tool] (out) {\code{updated\_*}\\file};
  \node[dec, right=of out] (rev) {human\\reviews};
  \node[box, below=11mm of rev] (imp) {human re-imports};
  \node[store, above=10mm of tool] (stock) {internal\\stock list};

  \draw[flow] (store) -- (exp);
  \draw[flow] (exp) -- (csv);
  \draw[flow] (csv) -- (tool);
  \draw[dflow] (stock) -- (tool);
  \draw[flow] (tool) -- (out);
  \draw[flow] (out) -- (rev);
  \draw[flow] (rev) -- (imp);
  \draw[flow] (imp) -| (store);
\end{tikzpicture}}
\caption{The round trip. The tools are one box; the human is four, and one of them
is a review gate no catalogue change can skip.}
\end{figure}

The constraint costs and pays. It pays because a wrong result is a wrong
spreadsheet rather than a wrong live storefront. It costs because without the API
there are no stable product IDs to join on, only whatever text the export happens
to contain. That is the whole engineering problem.

\section{The three tools}

\begin{center}
\begin{tabular}{@{}p{4.5cm}p{4.0cm}p{5.5cm}@{}}
\toprule
\textbf{Tool} & \textbf{Does} & \textbf{Notes} \\
\midrule
\code{quantity\_pricing\_updater} & Merges warehouse quantity and average cost into a Shopify export & CLI. 69 lines. \textbf{Verified working.} \\
\addlinespace
\code{dimensions\_extractor} & Pulls bore / outside diameter / width out of HTML descriptions into columns & CLI. 83 lines. \textbf{Verified working}, with a row-dropping defect. \\
\addlinespace
\code{stock\_updater\_gui} & tkinter app: both of the above, plus variant detection, behind three buttons & 354 lines. The client's deliverable. \textbf{Carries a serious bug.} \\
\bottomrule
\end{tabular}
\end{center}

Both CLI tools were shipped to the client as standalone Windows executables via
PyInstaller, so no Python install was needed.

\section{The central idea: matching without a key}

\begin{keyidea}{Two tables, no shared column}
The stock list has no Shopify ID. The Shopify export has no stock item number. The
only thing in both is the bearing designation as text, e.g. \code{6205}. And it is
not unique: \code{6205} appears once per brand, and only the \code{HI-TEC} row is
the one to trust. \textbf{Filter by brand first, then join.} That single decision
is what the project turns on.
\end{keyidea}

The join key then depends on the product's shape. Shopify gives every product at
least one variant, labelling a variant-less product's row
\code{Option1 Value = Default Title}. So:

\begin{itemize}
  \item \code{Default Title} rows: join on \code{Title}.
  \item Variant rows: join on \code{Option1 Value}, because Shopify writes the
  \code{Title} only on the first row of each product group and leaves it empty on
  the rest.
\end{itemize}

\begin{figure}[H]
\centering
\resizebox{0.92\textwidth}{!}{
\begin{tikzpicture}[node distance=7mm and 13mm]
  \node[store] (sl) {stock list};
  \node[box, right=of sl] (filt) {filter\\\code{Brand=='HI-TEC'}};
  \node[store, right=of filt] (hitec) {trusted rows};
  \node[store, below=15mm of sl] (ss) {Shopify export};
  \node[dec, right=of ss] (split) {\code{Default}\\\code{Title}?};
  \node[box, above right=1mm and 11mm of split] (defl) {join on \code{Title}};
  \node[box, below right=1mm and 11mm of split] (nond) {join on\\\code{Option1 Value}};
  \node[box, right=34mm of split] (cat) {\code{concat}};
  \node[store, right=of cat] (out) {+ \code{Updated Quantity}\\+ \code{Updated Cost}};
  \draw[flow] (sl) -- (filt); \draw[flow] (filt) -- (hitec);
  \draw[flow] (ss) -- (split);
  \draw[flow] (split) -- node[lbl,above,sloped]{yes} (defl);
  \draw[flow] (split) -- node[lbl,below,sloped]{no} (nond);
  \draw[flow] (defl) -| (cat); \draw[flow] (nond) -| (cat);
  \draw[flow] (cat) -- (out);
  \draw[dflow] (hitec) -- (defl);
\end{tikzpicture}}
\caption{Brand filter, then a split on the \code{Default Title} sentinel to choose
the join key.}
\end{figure}

\textbf{Verified by running it.} The shipped sample deliberately includes trap
rows: \code{6205} and \code{6304 2RS} each have a competing \code{OTHER} brand row
with different figures. The tool takes the HI-TEC figures for both, and matches
both variant rows through \code{Option1 Value} including the one with a blank
\code{Title}. The README's main claim about this tool is true.

\section{The second idea: data locked in HTML}

Bearing dimensions existed only as an HTML table inside each product's description,
so the storefront could not filter on them. The extractor lowercases the
description, strips every tag with the regex \code{<.*?>}, splits on newlines,
drops blanks, and then relies on one assumption: \textbf{a label line is
immediately followed by its value line}, so \code{['bore diameter', '25']} are
adjacent. Find the label, take the next line.

It works, and it recovered all three dimensions from every sample row that had a
description. It is also brittle by construction: it rides on the store's exact
description template and would silently miss dimensions on any other layout. The
README says this plainly and is right to.

One detail worth knowing: a \code{total width} fallback exists for products that
label the field differently. That line is domain knowledge learned from the real
catalogue, not something anyone would write from first principles.

\section{Findings}

\begin{honest}{The GUI's quantity matcher matches nothing on the shipped sample}
The GUI is not a wrapper around the two CLI tools; it is a copy of them, and the
copies have drifted. The GUI has one line its CLI sibling lacks:

\begin{lstlisting}[language=Python]
stock_list['Item'] = stock_list['Item'].apply(convert_int_float)
\end{lstlisting}

It is meant to normalise types. It does two damaging things. \code{'6205'} becomes
the float \code{6205.0}, which never compares equal to the string \code{'6205'} in
the Shopify \code{Title} column. And \code{'22210 K'} cannot parse as a float, so
the function falls through to a bare \code{pass}, returns \code{None}, and pandas
stores \code{nan}: every alphanumeric designation, meaning every variant, is
destroyed.

Running the GUI's exact matching logic against the sample: \textbf{0 of 4}
default-title matches, \textbf{0 of 2} variant matches. The CLI matches all six.
The GUI would write zeros for every product, report ``Quantity and Pricing has been
updated!'' in green, and save the file.

\textbf{Scope, stated precisely.} What is verified is the failure on the sample in
this repository. Whether it failed on the client's real 2023 files is
\textbf{unresolvable from this repo}, because it depends on the dtypes pandas
inferred from data that is no longer here. If the real columns were purely numeric
the line was harmless; if either contained one alphanumeric designation, and the
sample says the real catalogue was full of them, the match collapses exactly as it
does here. The second is far more likely.

The honest position, and a strong one: the CLI implementation is correct and
verifiably so; the GUI copy carries a well-understood dtype bug; and that bug is
the direct, predictable consequence of the code duplication the README already
names as the project's main design regret. The fix is two lines.
\end{honest}

\begin{honest}{The ``11,000-line Shopify export'' is not in this repository}
The README says, correctly and in the past tense, that the original repo held the
client's real data (about 1,600 stock rows and an 11,000-line Shopify export) and
that it was removed. Checked: what is actually here is a fabricated sample of
\textbf{8 stock rows} and \textbf{6 product rows} (117 physical lines, since the
HTML descriptions span many lines each).

So those figures rest entirely on the owner's recollection. Git history confirms
the real data was \textbf{never committed}, so it cannot be recovered from history
either. Good for confidentiality; it also means the numbers are unverifiable by
anyone without the original client files. Not an accusation that they are wrong, a
note that they are the one class of claim here that cannot be defended by pointing
at code. The sample demonstrates correctness, at six rows. It does not demonstrate
scale.
\end{honest}

\begin{honest}{The dimensions extractor deletes rows from its own output}
\code{df = df.dropna(subset=['Body (HTML)'])} removes every row without a
description. Since Shopify writes product-level fields only on the first row of a
product group, that is most variant rows. Verified: 6 rows in, 5 rows out; the
\code{22210 E} variant vanished. The tool's stated contract is ``adds three
columns''; its behaviour is ``adds three columns and silently drops rows''. The
one-line fix is to filter the loop's iteration source rather than the frame.
\end{honest}

\begin{honest}{Smaller, but real}
\begin{itemize}
  \item \code{save\_data} uses a full \code{read\_excel} parse as an existence
  check, and catches only \code{FileNotFoundError}. The previous output being open
  in Excel raises \code{PermissionError} instead, which is uncaught. For a client
  who lives in Excel, that is the likeliest real-world failure.
  \item \code{open\_file}'s retry loop cannot retry: both \code{if} branches
  return. Cancelling the dialog yields an empty path and an unhandled exception.
  \item Nothing writes to the status widget on failure, only on success.
  \item Matching is $O(n \times m)$: the stock list is re-filtered and re-listed on
  every Shopify row. One \code{merge} would be linear and shorter.
  \item The \code{Seal Type} mapping is three-quarters a no-op: only \code{2rs} is
  translated, the rest map to themselves.
  \item An unmatched product gets \code{Updated Quantity = 0}, indistinguishable
  from genuinely out of stock, in a column that feeds a live storefront.
  \item Debug \code{print}s throughout, in the version whose likely packaging
  (\code{console=False}, \textbf{inferred}: the GUI's \file{.spec} is absent)
  discards stdout entirely.
  \item No tests, though the sample data is already an adversarial fixture.
\end{itemize}
\end{honest}

\begin{keyidea}{Security: clean, by construction}
No credentials exist in this project because none can: no API client, no tokens, no
shop domain, no network calls. \file{.env.example} declares that nothing is needed.
A scan of the tree found nothing, and git history was checked: the client's real
data was never committed, so there is nothing to scrub. The sample is fabricated
and the README says so.
\end{keyidea}

\section{What is genuinely good}

\begin{itemize}
  \item \textbf{The brand-filter insight.} Seeing that the join is ambiguous and
  that \code{Brand} is the disambiguator is the thinking the project turns on. It
  removed a class of wrong match no fuzzy title matching would have fixed.
  \item \textbf{Reading Shopify's export conventions correctly}, including the
  \code{Default Title} sentinel and the empty \code{Title} on variant rows.
  \item \textbf{Domain knowledge that came from the data}, not from assumptions:
  the \code{total width} fallback, and a series-2xxxx rule in the variant detector
  that deliberately overrides the cage lookup table.
  \item \textbf{Delivery treated as a real requirement.} The GUI, the executables,
  the file pickers, and a non-overwrite save that bumps a counter exist because the
  user was a person rather than a terminal.
  \item \textbf{An honest scrub.} Replacing client data with a fabricated fixture
  that keeps the same column layout, the same HTML structure, and deliberate brand
  collisions, so the tools still run and can be verified, is better than most
  portfolio scrubs manage.
\end{itemize}

\section{What was verified, and what could not be}

\begin{center}
\begin{tabular}{@{}ll@{}}
\toprule
CLI merge picks HI-TEC rows, matches variants & \textbf{Verified by running it} \\
CLI extractor recovers all three dimensions & \textbf{Verified by running it} \\
Extractor drops description-less rows & \textbf{Verified by running it} \\
GUI matcher matches nothing on the sample & \textbf{Verified by running its logic} \\
GUI window matches the shipped screenshot & \textbf{Verified} \\
No credentials in tree or history & \textbf{Verified} \\
GUI's real \file{.spec} used \code{console=False} & \textbf{Inferred}, not in repo \\
Shopify import ignores absent rows & \textbf{Inferred}, not tested \\
GUI bug's effect on the real 2023 data & \textbf{Unresolvable} from this repo \\
1,600 stock rows / 11,000-line export & \textbf{Unverifiable}, never committed \\
End-to-end live round trip (2023) & \textbf{Not re-verified}, as README says \\
\bottomrule
\end{tabular}
\end{center}

\begin{keyidea}{The fair summary}
Three small tools that solved a real and genuinely awkward problem under a
constraint that forced text matching on data with no shared key. The central idea
is sound and verifiably works in the CLI tools. The delivered GUI is a copy-paste
fork that carries a serious silent bug, and the absence of both tests and a shared
module is exactly why that bug could survive unnoticed. The author has already
named most of this project's weaknesses accurately in his own README; the GUI
matcher is the one he had not found, and it is the most interesting thing here.
\end{keyidea}

\end{document}
